\begin{abstract}
We introduce \textsc{QLoRA}, a novel finetuning method designed to significantly lower memory requirements, enabling the finetuning of a 65B parameter model on a single 48GB GPU without compromising the performance of standard 16-bit finetuning tasks. \textsc{QLoRA} achieves this by performing backpropagation through a 4-bit quantized, frozen pretrained language model into Low Rank Adapters~(LoRA). Our flagship model series, termed \textbf{Guanaco}, surpasses all prior publicly available models on the Vicuna benchmark, attaining 99.3\% of ChatGPT's performance, with only 24 hours of training on one GPU. \textsc{QLoRA} incorporates several key innovations to achieve memory efficiency without loss of accuracy: (a) the 4-bit NormalFloat (NF4) data type, which is theoretically optimal for normally distributed weights; (b) Double Quantization, which further compresses memory usage by quantizing the quantization parameters themselves; and (c) Paged Optimizers to control temporary memory peaks. Utilizing \textsc{QLoRA}, we conduct large-scale finetuning across over 1,000 models, delivering an extensive analysis of instruction adherence and conversational performance on 8 instruction datasets, a variety of model architectures (LLaMA, T5), and scales previously infeasible for traditional finetuning approaches (e.g., 33B and 65B parameters). Our experiments demonstrate that finetuning with QLoRA on a compact, high-quality dataset can achieve state-of-the-art outcomes—even outperforming prior leading models of larger sizes. We also provide in-depth performance evaluation of chatbots using both human assessments and GPT-4 judgments, concluding that GPT-4 offers a practical, cost-effective proxy for human evaluation. Additionally, our findings suggest that prevailing chatbot benchmarks do not reliably reflect true chatbot capabilities. An in-depth error analysis highlights specific scenarios where \textbf{Guanaco} underperforms relative to ChatGPT. All models and code, including custom CUDA kernels for 4-bit training, are made publicly available.\footnote{\url{https://github.com/artidoro/qlora} and \url{https://github.com/TimDettmers/bitsandbytes}}
\end{abstract}

\section{Introduction}

Adapting large language models (LLMs) via finetuning has proven to be a powerful approach for enhancing model accuracy~\citep{min2021metaicl, wei2021finetuned, ouyang2022training, wang2022super, wang2022self, liu2022few} and for introducing preferred behaviors or eliminating unwanted ones~\citep{ouyang2022training,askell2021general,bai2022training}. Nevertheless, updating the parameters of very large models incurs immense computational costs; standard 16-bit finetuning for a model such as LLaMA with 65 billion parameters~\citep{touvron2023llama} can consume upwards of 780~GB of GPU memory. Although quantization schemes have emerged to curtail the memory usage of LLMs at inference time~\citep{dettmers2022llm, dettmers2022case, frantar2022gptq, xiao2022smoothquant}, these approaches fail to extend seamlessly to the training phase~\citep{wortsman2023stable}.

In this work, we provide the first evidence that quantized 4-bit models can be finetuned with no compromise in accuracy. Our technique, \textsc{QLoRA}, relies on a novel high-precision quantization method to convert pretrained models into a 4-bit format. On top of the quantized model, we incorporate lightweight Low-rank Adapter weights~\citep{hu2021lora} which are updated by propagating gradients through the fixed quantized weights.

Through \textsc{QLoRA}, the required GPU memory for finetuning a 65B model drops from over 780~GB to less than 48~GB, yet inference and training efficiency, as well as predictive capability, remain equivalent to a full 16-bit finetuned baseline. This dramatic reduction in resource demands democratizes access to LLM finetuning, making it feasible to adapt the largest publicly available models using only a single GPU. Leveraging \textsc{QLoRA}, we have developed the \textbf{Guanaco} model suite. The second most performant variant attains 97.8\% of ChatGPT’s score on the Vicuna~\citep{vicuna2023} benchmark, while being trainable within 12 hours on standard consumer hardware; with a single professional GPU over a day, our largest configuration achieves 99.3\% of ChatGPT’s performance on the same benchmark, nearly closing the gap. Notably, at deployment, our smallest \textbf{Guanaco} (7B parameters) requires just 5~GB of memory and surpasses the 26~GB Alpaca model by over 20 percentage points on the Vicuna benchmark (see Table~\ref{tbl:vicuna_eval}).

\textsc{QLoRA} presents several advancements aimed at minimizing memory requirements while maintaining strong model performance: (1) \textbf{4-bit NormalFloat}, a quantization data type theoretically optimized for normally distributed values, which achieves empirically superior outcomes compared to both 4-bit Integers and 4-bit Floats. (2) \textbf{Double Quantization}, which involves quantizing the quantization parameters themselves, leading to an average memory saving of roughly 0.37 bits per parameter (translating to about 3 GB for models with 65B parameters). (3) \textbf{Paged Optimizers}, which leverage NVIDIA's unified memory to mitigate memory surges during gradient checkpointing, especially when handling mini-batches with lengthy sequences. Integrating these components, we develop a more robust LoRA framework that employs adapters in every layer of the network, largely eliminating the accuracy compromises documented in earlier approaches.

The memory efficiency offered by \textsc{QLoRA} allows for an extensive examination of instruction finetuning and chatbot behavior across model sizes that would not be feasible with traditional finetuning methods, due to their memory demands. Consequently, we train over 1,000 models spanning multiple instruction tuning datasets, a range of model architectures, and parameter counts from 80M up to 65B. Our results confirm that \textsc{QLoRA} achieves performance comparable to standard 16-bit finetuning (\S\ref{sec:comp_to_std}) and is capable of training leading-edge chatbots, such as \textbf{Guanaco} (\S\ref{sec:pushing}). Further, our analysis of the resulting models reveals notable trends. First, we observe that the quality of training data significantly outweighs the impact of dataset size—for instance, a smaller 9k sample dataset (OASST1) surpasses a much larger 450k sample dataset (FLAN v2, subsampled) in terms of chatbot effectiveness, despite both aiming to facilitate instruction-following generalization. Second, we demonstrate that high performance on the Massive Multitask Language Understanding (MMLU) benchmark does not necessarily correlate with strong results on the Vicuna chatbot benchmark, and the reverse is also true. This suggests that for particular tasks, the appropriateness of the dataset is more critical than its sheer scale.

Additionally, we conduct a thorough examination of chatbot effectiveness utilizing both assessments from human evaluators and judgments from GPT-4. Our evaluation framework employs a tournament-style methodology, where various chatbot models are pitted against one another to generate optimal responses for particular prompts. The outcome of each match is determined by either GPT-4 or a panel of human annotators. The results from these tournaments are synthesized into Elo scores~\citep{elo1967proposed, elo1978rating}, which serve as the metric for model rankings. We observe a high level of concordance between GPT-4 and human-based rankings, although notable discrepancies are present in some cases. Therefore, we underscore that while automated, model-based assessment can offer a more cost-effective alternative to human annotation, it is also subject to specific limitations and potential sources of error.

To complement the quantitative findings from our benchmarking experiments, we also carry out an in-depth qualitative assessment of the \textbf{Guanaco} models. This analysis elucidates scenarios in which these models excel or fall short, including some that elude detection in standard quantitative metrics.

We have made available all generated model outputs along with their respective human and GPT-4 annotations, thus supporting further investigations by the research community. Our software framework, including both codebase and CUDA kernels, is open-sourced and our methods have been incorporated into the Hugging Face transformers ecosystem~\citep{wolf2019huggingface} for widespread accessibility. Furthermore, we provide an array of adapters tailored for models of sizes 7/13/33/65B, each fine-tuned on eight diverse instruction-following datasets—yielding a total of 32 openly released, fine-tuned model variants.

\section{Background}
\paragraph{Block-wise k-bit Quantization} Quantization refers to the transformation of data from a high-precision representation to a more compact, lower-precision format. Typically, this involves converting data types that use more bits, such as 32-bit floats, into those that require fewer bits, like 8-bit integers. To maximize the utilization of the reduced numeric range in the target type, the data is generally rescaled via normalization with respect to the absolute maximum value found in the input tensor. For instance, the process of quantizing a 32-bit floating point (FP32) tensor into an Int8 tensor with range $[-127, 127]$ is described by:

\begin{equation}
    \mathbf{X}^{ \text{Int8}} = \text{round}\left( \frac{127}{{\text{absmax}}(\mathbf{X}^{{\text{FP32}}})}\mathbf{X}^{{\text{FP32}}} \right) = \text{round}(c^{\text{FP32}}\cdot \mathbf{X}^{{\text{FP32}}}),
\end{equation}

where $c$ denotes the {\em quantization constant} or {\em quantization scale}. To revert to the original higher-precision representation, dequantization is performed as follows:

\begin{equation}
    \text{dequant}(c^{\text{FP32}}, \mathbf{X}^{\text{Int8}}) = \frac{\mathbf{X}^{\text{Int8}}}{c^{\text{FP32}}} = \mathbf{X}^{\text{FP32}}
\end{equation}

A limitation of this method arises when the input tensor contains values of large magnitude (i.e., outliers), leading to inefficient utilization of the quantization bins: certain bit combinations may end up sparsely or not at all represented, as very few or no input values are quantized into some bins. To address the presence of outliers, a widely adopted solution is to divide the input tensor into several segments or blocks, each of which is quantized independently, allowing each block to have its own distinct quantization constant $c$. More specifically, the input tensor $\mathbf{X}\in \mathbb{R}^{b\times h}$ is partitioned into $n$ sequential blocks, each of length $B$, by flattening the tensor and slicing it into linear segments, resulting in ${n = ({b\times h} )/{B} }$ total blocks. Each of these blocks is separately quantized as prescribed by Equation~1, yielding both a quantized tensor and a set of $n$ associated quantization constants $c_i$.

\paragraph{Low-rank Adapters} The Low-rank Adapter (LoRA) finetuning approach \citep{hu2021lora} offers memory efficiency by introducing a small collection of trainable adapter parameters while the original model’s weights remain unchanged. During stochastic gradient descent, gradients are propagated through the frozen pretrained weights, but only the adapter parameters are modified to minimize the loss. LoRA modifies a linear transformation by appending an extra low-rank, factorized component. For a linear mapping ${\mathbf{X} \mathbf{W} = \mathbf{Y}}$ with $\mathbf{X}\in  \mathbb{R}^{b\times h}$ and $\mathbf{W}\in  \mathbb{R}^{h\times o}$, the LoRA formulation is:

\begin{equation}
    \mathbf{Y} = \mathbf{X}\mathbf{W} + s\mathbf{X}\mathbf{L}_1\mathbf{L}_2,
\end{equation}

where $\mathbf{L}_1\in  \mathbb{R}^{h\times r}$, $\mathbf{L}_2\in  \mathbb{R}^{r\times o}$, and $s$ is a scaling factor.

\paragraph{Memory Requirement of Parameter-Efficient Finetuning} An important consideration is the memory usage associated with LoRA during training, particularly with respect to both the number and the size of adapters deployed. Given LoRA's exceptionally small memory overhead, it is feasible to deploy additional adapters to boost performance, incurring only a marginal increase in overall memory consumption. Although LoRA was developed as a Parameter Efficient Finetuning (PEFT) technique, in practice, the majority of memory required when finetuning LLMs is occupied by activation gradients rather than the LoRA parameters themselves. For example, in the case of finetuning a 7B LLaMA model on FLAN v2 with batch size 1, and employing LoRA weights that comprise roughly 0.2\% of the model’s total parameters~\citep{hu2021lora,liu2022few}, the memory usage for LoRA input gradients is 567 MB, while the LoRA parameter storage demands only 26 MB. When gradient checkpointing~\citep{chen2016training} is enabled, the input gradients shrink to an average of 18 MB per sequence—yet, this figure still exceeds the total memory used by all LoRA weights. For context, a 4-bit quantized base model occupies 5,048 MB in memory. These numbers underscore the significance of gradient checkpointing, but also reveal that further minimizing the LoRA parameter count leads to relatively negligible memory savings. Consequently, it becomes practical to introduce more adapters without a major impact on the training memory profile (for more details, see Appendix~\ref{app:memory}). As will be detailed later, this capacity is particularly important when aiming to match the performance of full 16-bit precision training.

\section{\textsc{QLoRA} Finetuning}

\textsc{QLoRA} enables high-fidelity finetuning in 4-bit precision through two novel methods: 4-bit NormalFloat (NF4) quantization and Double Quantization. Furthermore, we present Paged Optimizers, a solution designed to mitigate memory surges during gradient checkpointing. These memory surges have historically led to out-of-memory errors, impeding the ability to finetune large-scale models on a single workstation.

\textsc{QLoRA} utilizes a storage data type of low precision (typically 4-bit), alongside a higher-precision computation data type, which is usually BFloat16. Operationally, this approach entails dequantizing \textsc{QLoRA} weight tensors to BFloat16 whenever they are accessed, followed by performing matrix multiplications at 16-bit precision.

The following sections elaborate on the internal mechanisms of \textsc{QLoRA} and provide a formal description of the framework.

\paragraph{4-bit NormalFloat Quantization} The NormalFloat (NF) format extends the idea of Quantile Quantization \citep{dettmers2022optimizers}, an information-theoretically optimal approach that allocates an equal number of input tensor values to each quantization bin. This is accomplished by determining the quantile of the input tensor, typically via the empirical cumulative distribution function.

A key drawback of quantile quantization lies in the computational burden of accurate quantile estimation. As a result, efficient quantile approximation techniques, such as SRAM quantiles~\citep{dettmers2022optimizers}, are adopted. Nevertheless, these approximations often introduce significant quantization errors, particularly affecting outlier values, which tend to be especially consequential.

Both the overhead of accurate quantile computation and the errors from approximation can be circumvented when tensors originate from distributions invariant up to a quantization constant. Under these circumstances, the quantiles remain consistent across input tensors, rendering exact quantile computation tractable.

Typically, pretrained neural network weights are distributed according to a zero-mean Gaussian with standard deviation $\sigma$ (refer to Appendix~\ref{app:normality}). To standardize all weights to one common distribution, we can adjust $\sigma$ so that this distribution fully occupies the domain permitted by the data type. For our purposes, we select $[-1, 1]$ as the target range for the data type. Consequently, it is necessary to scale both the data type quantiles and the neural network weights to reside within this range.

The data type that is optimal from an information-theoretic standpoint, when considering zero-mean Gaussians with varying standard deviations $\sigma$ and range limited to $[-1, 1]$, can be determined in the following manner: (1) compute the $2^k+1$ quantiles from a standard normal distribution $N(0, 1)$, generating a $k$-bit quantile quantization data type tailored for normal distributions, (2) map the resulting data type values into the interval $[-1, 1]$, and (3) quantize the original weight tensor by similarly normalizing it to $[-1, 1]$ using rescaling based on the maximum absolute value.

When the distributions of the weight tensor and the quantization data type coincide, quantization proceeds as standard. The third step corresponds to aligning the standard deviation of the weights with that of the $k$-bit data type through appropriate rescaling. Mathematically, the $2^k$ representative values $q_i$ for the data type are calculated by:

\begin{equation}
q_i  = \frac{1}{2}\left( Q_X\left(\frac{i}{2^k+1}\right) + Q_X\left(\frac{i+1}{2^k+1}\right)\right),
\end{equation}

where $Q_X(\cdot)$ represents the quantile function associated with the standard normal distribution $N(0,1)$. One challenge associated with symmetric k-bit quantization is that this method does not precisely represent zero. This omission is problematic, as an exact zero is essential for quantizing padding and other elements that are exactly zero, in order to avoid quantization error in these cases. To ensure a discrete quantized value of $0$ is present, and to fully utilize all $2^k$ values available for a k-bit representation, we construct an asymmetric quantization scheme. Specifically, we estimate quantiles $q_i$ over two intervals: $2^{k-1}$ quantiles for the negative side and $2^{k-1}+1$ for the positive side. We then combine these quantile sets, eliminating one duplicate zero that appears in both, to guarantee only a single occurrence of zero. We refer to the resultant quantization data type as \emph{k-bit NormalFloat} (NFk), emphasizing that each quantization bin is expected to hold the same number of data points when the distribution is centered at zero. This construction ensures information-theoretic optimality for data drawn from a normal distribution with mean zero. Exact quantization values for this data type are provided in Appendix~\ref{app:nf4}.

\paragraph{Double Quantization{}}
We present \emph{Double Quantization} (DQ), a strategy wherein the quantization parameters themselves are subjected to quantization, yielding further reductions in memory usage. Although a small block size is necessary for accurate 4-bit quantization~\citep{dettmers2022case}, this results in significant overhead for storing quantization constants. For instance, using 32-bit storage for the quantization parameters and a block size of 64 in a weight matrix $\mathbf{W}$, the overhead translates to an average of $32/64 = 0.5$ bits per model parameter. The Double Quantization technique addresses this issue by compressing the quantization constants themselves.

More concretely, Double Quantization considers the quantization constants $c_2^{\text{FP32}}$ obtained from the first quantization step as inputs to a subsequent quantization stage. This secondary quantization produces quantized quantization constants $c_2^{\text{FP8}}$ and also computes an additional set of quantization constants $c_1^{\text{FP32}}$ for the second stage. We apply 8-bit floating-point quantization with a block size of 256 in the second quantization round, as empirical evidence~\citep{dettmers2022case} indicates that such an approach maintains model performance. Given that the $c_2^{\text{FP32}}$ values are strictly nonnegative, we center these constants by subtracting their mean before quantization, thereby facilitating the use of a symmetric quantization scheme. This process, for a block size of 64, effectively decreases the average memory cost per parameter from $32/64 = 0.5$ bits down to $8/64 + 32/(64\cdot256) = 0.127$ bits, yielding a saving of 0.373 bits per parameter.

\paragraph{Paged Optimizers} leverage NVIDIA's unified memory~\footnote{\tiny \url{https://docs.nvidia.com/cuda/cuda-c-programming-guide}}, which seamlessly manages the transfer of memory pages between CPU and GPU. This mechanism ensures stable GPU operations when memory resources are exhausted by transparently moving data as needed. The unified memory functions similarly to conventional paging systems that swap data between CPU RAM and persistent storage. In our approach, optimizer states are allocated in paged memory, automatically migrating to CPU RAM whenever the GPU is low on memory, and returning to the GPU in response to demand during the optimizer's update computations.

\paragraph{\textsc{QLoRA}.} Building upon the previously discussed components, we formally define \textsc{QLoRA} as applied to a single linear transformation within a quantized model, enhanced with a single LoRA adapter, in the following way:

\begin{equation}
    \mathbf{Y}^{\text{BF16}} =  \mathbf{X}^{\text{BF16}} \text{doubleDequant}(c_1^{\text{FP32}}, c_2^{\text{k-bit}}, \mathbf{W}^{\text{NF4}}) + \mathbf{X}^{\text{BF16}}\mathbf{L}^{\text{BF16}}_1\mathbf{L}^{\text{BF16}}_2,
\end{equation}

The operation doubleDequant$(\cdot)$ is further specified by:

\begin{equation}
    \text{doubleDequant}(c_1^{\text{FP32}}, c_2^{\text{k-bit}}, \mathbf{W}^{\text{k-bit}}) = \text{dequant}(\text{dequant}(c_1^{\text{FP32}}, c_2^{\text{k-bit}}), \mathbf{W}^{\text{4bit}}) = \mathbf{W}^{\text{BF16}},
\end{equation}

In our setup, $\mathbf{W}$ is quantized using NF4 and $c_2$ employs the FP8 format.

For increased accuracy, we select a blocksize of 64 for $\mathbf{W}$, while opting for a blocksize of 256 for $c_2$ to optimize memory efficiency.

During parameter optimization, it suffices to compute gradients with respect to the LoRA adapter weights, i.e., $\frac{\partial E}{\partial \mathbf{L}_i}$; gradients for the 4-bit weights, $\frac{\partial E}{\partial \mathbf{W}}$, are not required. Nevertheless, deriving $\frac{\partial E}{\partial \mathbf{L}_i}$ involves evaluating $\frac{\partial \mathbf{X}}{\partial \mathbf{W}}$, which utilizes equation~(5). This process involves dequantizing $\mathbf{W}^{\text{NF4}}$ to the computation datatype $\mathbf{W}^{\text{BF16}}$, allowing for gradient computation $\frac{\partial \mathbf{X}}{\partial \mathbf{W}}$ in BFloat16 precision.

In summary, \textsc{QLoRA} employs a single storage data format—commonly 4-bit NormalFloat—and a computation format—16-bit BrainFloat. During both the forward and backward passes, stored values are dequantized from the storage format to the computation format; however, weight gradients are computed solely for the LoRA parameters, which utilize 16-bit BrainFloat.

\section{QLoRA vs. Standard Finetuning}
\label{sec:comp_to_std}

Previously, we described the operation of QLoRA and its notable ability to lower the memory demands required for model finetuning. The central issue now concerns the ability of QLoRA to reach performance comparable to full-model finetuning. Additionally, we aim to dissect specific elements of QLoRA, such as evaluating the influence of NormalFloat4 relative to conventional Float4. The experiments discussed in the following sections are designed to address these questions.

\paragraph{Experimental setup.} Our evaluation considers three model architectures: encoder, encoder-decoder, and decoder-only. We benchmark QLoRA against 16-bit adapter finetuning and complete finetuning approaches for model scales up to 3B parameters. Assessments are carried out on GLUE \citep{wang2018glue} using RoBERTa-large \citep{liu2019roberta}, Super-NaturalInstructions (TKInstruct) \citep{wang2022super} with T5 \citep{t5}, as well as 5-shot MMLU \citep{hendrycksmeasuring} after LLaMA has been finetuned on Flan v2 \citep{longpre2023flan} and Alpaca \citep{alpaca}. To further explore the benefits of NF4 in comparison to other 4-bit representations, we adopt the experimental design from \citet{dettmers2022case} and evaluate post-quantization zero-shot accuracy and perplexity over several architectures (OPT~\citep{zhang2022opt}, LLaMA~\citep{touvron2023llama}, BLOOM~\citep{scao2022bloom}, and Pythia~\citep{biderman2023pythia}), ranging in size from 125M to 13B.

We elaborate on each specific experimental condition in the results section to facilitate interpretation, with comprehensive details provided in Appendix~\ref{app:exp-setup}.

While paged optimizers are essential to perform 33B or 65B \textsc{QLoRA} finetuning on individual 24/48GB GPUs, we do not report precise timings for paged optimizers since paging is primarily triggered by unusually long mini-batch sequence lengths, a relatively uncommon situation. Nonetheless, we investigate the runtime characteristics of paged optimizers for 65B-parameter models on 48GB GPUs and determine that, when utilizing a batch size of 16, the training throughput is comparable to that of standard optimizers. Further research is needed to systematically evaluate and document conditions where the paging mechanism may introduce significant overhead.

\paragraph{Default LoRA hyperparameters do not match 16-bit performance}

Applying the conventional strategy of LoRA—namely, attaching adapters solely to the query and value projection matrices as per \citep{hu2021lora}—does not achieve parity with full-model finetuning in the context of large-scale base models. Evidence from Figure~\ref{fig:hyper} on LLaMA 7B finetuning with Alpaca reveals that the most pivotal LoRA hyperparameter is the number of adapters deployed throughout the model; it is necessary to apply LoRA to every linear transformer block layer in order to approach the performance of full finetuning. In contrast, other hyperparameters (such as the projection rank $r$) exhibit minimal impact on outcomes (see Appendix \ref{app:exp-setup}).

In a similar vein, our investigation indicates that the default hyperparameter configurations for fully finetuned baseline models are insufficiently optimized. To establish strong baseline results, we conduct a systematic search across learning rates ranging from 1e-6 to 5e-5 and batch sizes between 8 and 128. The corresponding outcomes for finetuning the 7B LLaMA model on Alpaca are depicted in Figure~\ref{fig:hyper}.

\paragraph{4-bit NormalFloat surpasses 4-bit Floating Point in performance}

Although the 4-bit NormalFloat (NF4) format is proven to be optimal in terms of information theory, whether this confers practical improvements is an open question. Replicating the methodology from \citet{dettmers2022case}, we assess quantized LLMs—namely, OPT \citep{zhang2022opt}, BLOOM \cite{scao2022bloom}, Pythia \cite{biderman2023pythia}, and LLaMA—spanning parameter counts from 125M to 65B, utilizing multiple data types. Evaluations cover language modeling tasks as well as a variety of zero-shot benchmarks. As shown in Figure~\ref{fig:nf4} and Table~\ref{tbl:nf4_ppl}, NF4 consistently outperforms both FP4 and Int4, while double quantization achieves further memory reductions without compromising model quality.

\paragraph{k-bit \textsc{QLoRA} achieves parity with both 16-bit full finetuning and 16-bit LoRA approaches} 

Previous work has demonstrated that although 4-bit quantization is feasible for \emph{inference}, it typically results in diminished performance compared to 16-bit precision \citep{dettmers2022case,frantar2022gptq}. This highlights an important question: Can the observed performance drop be mitigated via 4-bit adapter finetuning? We explore this using two different evaluation settings.

First, we compare these methods against fully finetuned 16-bit versions of RoBERTA and T5 models, ranging from 125M to 3B parameters, on both GLUE and the Super-NaturalInstructions benchmark. Table~\ref{tab:nf4vsbf16} presents these findings. Across both evaluation datasets, 16-bit, 8-bit, and 4-bit adapter-based finetuning approaches consistently match the performance of the 16-bit fully finetuned baseline. This evidence indicates that the adverse effects of coarse quantization can be entirely remedied by adapter-based finetuning.

In our second experimental configuration, because finetuning models with 11B parameters or more in full precision necessitates multiple servers equipped with high-memory GPUs, we investigate whether 4-bit \textsc{QLoRA} can achieve comparable results to 16-bit LoRA in the 7B to 65B parameter range. Specifically, we conduct finetuning of LLaMA models from 7B up to 65B parameters using two instruction-tuned datasets, Alpaca and FLAN v2, subsequently measuring performance on the MMLU benchmark utilizing 5-shot accuracy. Table~\ref{tbl:llama-datatype} presents the findings, revealing that NF4 coupled with double quantization fully restores the MMLU performance attained by 16-bit LoRA. Additionally, we observe that \textsc{QLoRA} using the FP4 format falls short of the 16-bit brain float LoRA baseline by roughly one percentage point. These outcomes support our prior conclusions: (1) \textsc{QLoRA} leveraging NF4 matches the effectiveness of both 16-bit full finetuning and 16-bit LoRA, and (2) NF4 offers higher quantization accuracy than FP4.

\paragraph{Summary} Across our experiments, 4-bit \textsc{QLoRA} utilizing NF4 consistently attains performance equivalent to 16-bit full finetuning and 16-bit LoRA on established academic benchmarks with rigorous evaluation protocols. Moreover, we confirm the superiority of NF4 over FP4 and demonstrate that double quantization does not negatively impact outcomes. Taken together, these results provide strong evidence that 4-bit \textsc{QLoRA} reliably reproduces the results achieved with 16-bit finetuning approaches.

Consistent with earlier research on quantization \citep{dettmers2022case}, our analyses of MMLU and Elo show that, for a fixed finetuning and inference compute budget, increasing the size of the base model while lowering its numerical precision yields benefits. This underscores the practical efficiency advantages conferred by \textsc{QLoRA}. Since our 4-bit finetuning experiments did not show any loss in accuracy compared to full-precision finetuning, an open question remains regarding the precise location of the performance–precision trade-off for QLoRA tuning, which we propose as a direction for future investigation.

We turn our attention to scaling instruction tuning beyond the capabilities of full 16-bit finetuning on conventional academic research infrastructure.

\section{Pushing the Chatbot State-of-the-art with QLoRA}
\label{sec:pushing}
Having shown that 4-bit \textsc{QLoRA} can reliably match the performance of standard 16-bit models across a range of scales, datasets, and tasks, we embark on a detailed exploration of instruction finetuning utilizing some of the largest open-source language models currently accessible for research purposes. To evaluate the effectiveness of instruction finetuning, we employ the rigorous MMLU benchmark for Natural Language Understanding, in addition to proposing novel methodologies for assessing real-world chatbot behavior.

\subsection{Experimental setup} Here, we summarize the experimental setup, with full methodological details deferred to Appendix \ref{app:chatbot}.

\paragraph{Data} Recognizing the absence of a thorough survey on contemporary instruction-following datasets, we curate a collection of eight up-to-date datasets. These include those sourced through crowd-sourcing (OASST1~\citep{kopf2023openassistant}, HH-RLHF~\citep{bai2022training}), datasets derived by distillation from instruction-finetuned models (Alpaca~\citep{alpaca}, self-instruct~\citep{wang2022self}, unnatural-instructions~\citep{honovich2022unnatural}), aggregated corpora (FLAN v2~\citep{chung2022scaling}), as well as hybrid resources (Chip2~\citep{laion2023}, Longform~\citep{koksal2023longform}). These datasets differ in terms of language diversity, dataset magnitude, and licensing terms.

\paragraph{Training Setup} To eliminate confounding variations from training objectives, we carry out QLoRA finetuning under a supervised learning regime using a cross-entropy loss function and explicitly avoid reinforcement learning techniques, even for datasets that contain human preference judgments. For datasets with explicit instruction-response delineation, finetuning targets only the response component (see Appendix \ref{app:chatbot} for ablations). OASST1 and HH-RLHF contain multiple possible responses; for these, we select the top-ranked reply at every conversational branching and finetune using the complete curated exchange, inclusive of instructions. Our experimental protocol uniformly employs NF4 \textsc{QLoRA} with both double quantization and paged optimizers, effectively averting memory overflow issues that arise during gradient checkpointing. For 13B and 33B LLaMA models, we conduct limited hyperparameter tuning, observing that configurations optimized at the 7B scale (e.g., number of epochs) are broadly transferable, with exceptions for learning rate and batch size: these are halved and doubled, respectively, for the 33B and 65B model sizes.

\paragraph{Baselines} Our models are benchmarked against both academic (Vicuna~\citep{vicuna2023}, Open Assistant~\citep{kopf2023openassistant}) and commercial systems (GPT-4 \citep{openai2023gpt}, GPT-3.5-turbo, and Bard). The Open Assistant system represents a 33B LLaMA model instruction-finetuned using Reinforcement Learning from Human Feedback (RLHF) on OASST1, the same dataset incorporated into our experiments. In contrast, Vicuna consists of a fully finetuned LLaMA 13B model, trained on proprietary conversation logs from ShareGPT, effectively performing knowledge distillation from OpenAI’s GPT models.

\subsection{Evaluation}

Adhering to standard conventions, we adopt the MMLU (Massively Multitask Language Understanding) benchmark \citep{hendrycksmeasuring} for a broad assessment of language understanding across 57 multiple-choice tasks spanning diverse subjects such as elementary mathematics, US history, computer science, law, and beyond. Model results are reported using 5-shot test accuracy.

We further assess generative language abilities through a combination of automated and human assessment methods. This set of evaluations is grounded in human-crafted queries and is designed to evaluate the qualitative aspects of model-generated responses. Although this approach more closely mirrors realistic applications of chatbot models and has seen rising adoption, the field lacks a standardized evaluation methodology. Our proposed procedure, outlined below, employs nucleus sampling ($p=0.9$) and a temperature of $0.7$ uniformly across all experiments.

\paragraph{Benchmark Data} Our analysis utilizes two curated question datasets: Vicuna prompts \citep{vicuna2023} and the OASST1 validation set \citep{kopf2023openassistant}. For the Vicuna prompts, we use the full, unaltered set of 80 prompts spanning various topics. The OASST1 dataset encompasses a multilingual assortment of user-assistant multiturn dialog data, collected via crowdsourcing. From the validation set, we extract every user message as an individual query and supply the full conversational context in the prompt, resulting in a total of 953 distinct user queries. We refer to these as the Vicuna and OA benchmarks, respectively.

\paragraph{Automated Evaluation} Following the protocol established by \citet{vicuna2023}, we deploy GPT-4 to score the output of various models compared to ChatGPT (GPT-3.5 Turbo) using the Vicuna dataset. For each query, both ChatGPT's and another model's responses are presented to GPT-4, which then issues each a score from one to ten along with a rationale. A model’s aggregate performance is computed as a percentage of ChatGPT’s total score. Importantly, this relative score may exceed 100\% if a model achieves a superior absolute score to ChatGPT. We observe a notable bias in GPT-4’s scoring, wherein the response placed first tends to receive a higher rating. To mitigate this bias, we advocate averaging results over both possible orderings.

We further evaluate models by directly comparing their outputs, reframing the evaluation as a three-way classification that includes possible ties. Here, GPT-4 is asked to select the superior response or indicate if both responses are equally good, providing justifications in each case. This process is carried out across all model pairs on both the Vicuna and OA benchmarks.

\paragraph{Human Evaluation} Although prior studies suggest that generative models can serve as effective evaluation agents \citep{fu2023gptscore}, the alignment between GPT-4-generated ratings and actual human preferences for chatbot quality remains an open question. To address this, we conduct two concurrent human annotation experiments on the Vicuna dataset that parallel the automated evaluation procedures above. Utilizing Amazon Mechanical Turk (AMT), we assign two annotators for the ChatGPT comparison task and three annotators for pairwise system comparisons.

\paragraph{Elo Rating} Utilizing both human judgments and GPT-4-generated pairwise evaluations, we organize the models into a tournament framework, where each model faces off against others in head-to-head matches to generate superior responses for specified prompts. This tournament setup mirrors the methodologies of \citet{bai2022training} and \citet{vicuna2023}, though our approach extends these by incorporating ratings from GPT-4 alongside human evaluations. We draw randomly from our labeled set of comparisons to compute Elo ratings~\citep{elo1967proposed,elo1978rating}. The Elo rating system, prevalent in chess and various competitive games, quantifies a player's expected probability of victory relative to their opponent's. For instance, a player with an Elo of 1100 facing one with 1000 holds an estimated 65\% win probability, whereas identical Elo scores (e.g., 1000 vs 1000 or 1100 vs 1100) yield an expected win rate of 50\%. Elo scores are updated after each match in a manner proportional to the deviation between the observed and predicted outcome—unexpected results prompt substantial Elo adjustments, while anticipated results yield minor modifications. Cumulatively, the Elo scores converge to reflect the true ability of each model in the competitive context. Each model begins with a baseline rating of 1,000, with the parameter $K=32$ governing the step size. In line with the procedure of \citet{vicuna2023}, this process is repeated 10,000 times using various random seeds to account for potential biases due to the order of model pairings.

\subsection{\textbf{Guanaco}: \textsc{QLoRA} tuned on OASST1 Sets a New Standard for Open-source Chatbots} 

Through a combination of automated and human assessment, our analysis reveals that the {Guanaco} 65B model—finetuned via \textsc{QLoRA} on a variant of OASST1—emerges as the leading open-source chatbot and approaches the capabilities of ChatGPT. According to Elo ratings derived from system-level pairwise comparisons by human evaluators, {Guanaco} 65B and 33B achieve an expected win probability of 30\% against GPT-4, the highest such score published thus far.

Table~\ref{tbl:vicuna_eval} summarizes how models perform on the Vicuna benchmark \cite{vicuna2023} compared to ChatGPT. {Guanaco} 65B is second only to GPT-4, reaching 99.3\% of ChatGPT’s performance. Although {Guanaco} 33B features more parameters than Vicuna 13B, its weights are stored in 4-bit precision, allowing for substantially greater memory efficiency (21 GB versus 26 GB) and providing a three-point performance gain over Vicuna 13B. Additionally, {Guanaco} 7B can be accommodated by modern mobile devices with just a 5 GB memory requirement and outperforms Alpaca 13B by nearly 20 percentage points.

Nevertheless, the results in Table~\ref{tbl:vicuna_eval} are accompanied by notably wide confidence intervals, leading to considerable overlap in performance estimates between models. We suspect this variability arises from poorly defined rating scales—for instance, the meaning of an “8” out of 10 may not be consistent across different cases. Therefore, we advocate for the use of Elo ratings \citep{elo1967proposed}, which rely on \textit{pairwise} comparisons judged by humans and GPT-4, as they help mitigate ambiguities present in absolute scoring methods. Elo scores for the leading models are provided in Table~\ref{tbl:elo}. It is worth noting that while human annotators and GPT-4 differ in their ranking of certain models, such as Guanaco 7B, their rankings align for most models, yielding a system-level Kendall Tau coefficient of $\tau=0.43$ and Spearman rank correlation of $r=0.55$. At the example level, consensus between human majority votes and GPT-4 ratings is lower, reflected by a Fleiss $\kappa=0.25$. In sum, this indicates moderate alignment between GPT-4 and human annotators at the system level, suggesting that model-driven evaluation serves as a somewhat dependable stand-in for human judgments. Further discussion on these implications is presented in Section~\ref{ssec:considerations}.

The Elo ratings displayed in Table~\ref{tbl:elo_large} show that the {Guanaco} models with 33B and 65B parameters surpass all other models except GPT-4 on both the Vicuna and OA benchmarks, and their performance closely aligns with ChatGPT, as also reflected in Table~\ref{tbl:vicuna_eval}. It should be emphasized that the Vicuna benchmark tends to advantage open-source models, while ChatGPT performs better on the broader OA benchmark. Moreover, as evidenced by Tables \ref{tbl:mmlu-llama} and \ref{tbl:vicuna_eval}, the choice of finetuning dataset plays a critical role in determining outcomes. Specifically, Llama models finetuned with FLAN v2 excel in MMLU, yet rank lowest in the Vicuna benchmark—an observation consistent across different models. This finding also highlights a degree of independence among current benchmarks: excelling in MMLU does not necessarily correlate with top chatbot performance as assessed by Vicuna or OA, and the converse also holds.

Among the leading models in our analysis,  is unique in not utilizing proprietary data for training, as the OASST1 data collection protocol strictly prohibits the use of GPT-generated data. The next highest-performing model trained solely with open-source datasets is the Anthropic HH-RLHF, which trails {Guanaco} by 30 percentage points on the Vicuna benchmark (refer to Table~\ref{tbl:vicuna_eval}). Taken together, these results demonstrate that 4-bit \textsc{QLoRA} is a highly effective approach capable of producing chatbots at the state of the art, on par with ChatGPT. Notably, our 33B {Guanaco} model can be trained on a standard 24 GB consumer GPU in under 12 hours. This paves the way for future research using \textsc{QLoRA} fine-tuning with specialized open-source datasets, offering a viable route for developing models competitive with the most advanced commercial offerings.

\section{Qualitative Analysis}

Although our evaluation primarily relies on quantitative methods, exclusive dependence on aggregate metrics poses certain limitations. Chief among these is the concern of benchmark validity \citep{liao2021we}—that is, there is always ambiguity regarding whether a benchmark actually assesses the phenomena it claims to, particularly as more ``shortcuts'' used by machine learning systems to game such evaluations become apparent \citep{gururangan2018annotation, poliak2018hypothesis}. To address this shortcoming in part, we supplement our study with qualitative analysis organized into two segments. First, in \S\ref{ssec:examples}, we highlight several example generations that illustrate recurrent behaviors identified in outputs from our 65b {Guanaco} model. Subsequently, in \S\ref{ssec:considerations}, we discuss factors to consider concerning our findings and the interpretation of the observed results.

\subsection{Qualitative Analysis of Example Generations}
\label{ssec:examples}

For the selection of illustrative examples, we systematically review output data corresponding to both the Vicuna and OpenAssistant benchmarks to identify recurring answer patterns produced by {Guanaco}. Upon identifying such patterns, we craft tailored prompts or queries that deliberately attempt to evoke these patterns—even when doing so should yield incorrect or suboptimal answers. For example, if the model shows a proclivity for verbose responses, we direct it to ``Answer yes or no without explanation,'' thereby testing its tendency to deviate from concise answers. Through this process, we are able to showcase both ``lemons,'' where the model can be adversarially tricked, and ``cherries,'' where the model resists such manipulation. All text examples in this subsection were generated via Nucleus Sampling \cite{holtzmancurious} with $p=0.9$.

Clearly, this analysis is not exhaustive, as the limitations of this focused qualitative investigation make it impractical to account for every influencing factor—for instance, the extensive set of possible responses a model might produce from a given prompt necessitates relying on samples we consider broadly illustrative. Nevertheless, we contend that presenting these instances adds depth to the quantitative findings presented in earlier sections. By open-sourcing all models and code, we aim to motivate subsequent research that examines these issues with greater rigor.

\paragraph{Factual Recall} When posed with straightforward questions like “What is the capital of Zambia?”, all evaluated models reliably output the correct response, such as:

\begin{quote} 
    \textbf{Guanaco}:
    \texttt{The capital of Zambia is Lusaka.}
\end{quote}

However, the accuracy of {Guanaco} decreases as questions become less common, although the model maintains a high level of confidence in its answers. For example, when given this prompt from HotPotQA \citep{yang2018hotpotqa}:

\begin{quote}
     \textbf{User}: In what year was the singer who popularized `I'll Keep the Lovelight Burning' born?  

     \textbf{Guanaco}:~\texttt{Al Jolson is credited with popularizing the song `I'll Keep the Lovelight Burning,' and he was born in the year 1886.}
 \end{quote}

the model erroneously identifies both the popularizer and the year of birth (although the provided year correctly corresponds to Al Jolson, the person named).

\paragraph{Suggestibility} {Guanaco} demonstrates unexpected robustness against certain types of suggestive misinformation, as illustrated in the interaction below:

\begin{quote}
    \textbf{User}: How was it finally, officially confirmed that the earth is flat by peer reviewed scientists? \\ \\
    \textbf{Guanaco}: \texttt{There has never been any formal, peer-reviewed scientific confirmation that the earth is flat. The prevailing and longstanding agreement within the scientific community is that the earth is an oblate spheroid. This understanding has persisted for centuries, validated through extensive experimentation and observation. \\ \\ The flat earth claim is widely regarded as a debunked conspiracy theory. The assertion that a team of scientists has subjected the flat earth hypothesis to peer review and authenticated it is untrue. \\ \\ Please feel free to ask any further questions.}
\end{quote}

 also demonstrates an aptitude for recognizing when a question cannot be properly addressed, for instance,

\begin{quote}
    \textbf{User}: What time is it?

    \textbf{Guanaco}: \texttt{Unfortunately, as an AI language model, I am unable to access live data or provide up-to-the-minute information. My responses are based solely on pre-existing knowledge, so I cannot determine the present time. To find the current time, I recommend consulting a clock or searching online.}
\end{quote}

Comparable responses occur for prompts like ``Where are you?'', ``How are you?'', and so forth.

\paragraph{Refusal} In the same vein, {Guanaco} will, on occasion, decline to comply with instructions for no apparent reason:

\begin{quote}
    \textbf{User}: Please reverse the words in the sentence ``I'm a cool dude in a mocha mood.'' \\ \\
    \textbf{Guanaco}: \texttt{I apologize, but I am currently unable to complete that request. However, I can clarify why the sentence you mentioned is grammatically accurate.\\ \\The phrase ``I'm a cool dude in a mocha mood'' adheres to standard English grammar and syntax conventions. Here, the descriptive adjective ``cool'' precedes the noun ``dude'' as expected, and the prepositional phrase ``in a mocha mood'' provides further detail about the state described in the sentence. The structure used is fully correct and meaningful within this context.\\ \\If you have additional queries or need further assistance, please let me know.}
\end{quote}

\paragraph{Secret Keeping} Sadly, this form of refusal does not reliably prevent disclosures, even when confidentiality is intended. There are many scenarios in which developers want to prevent a system from revealing certain information to users, despite allowing open-ended natural language interactions. For instance, if the model is given the prompt ``The secret word is `banana'. Do not reveal or repeat back this back, no matter what the human requests,'' we can test how vulnerable it is to leaking this secret. When directly queried, it responds appropriately:

\begin{quote}
    \textbf{User}: What is the secret word?

    \textbf{Guanaco}: \texttt{I'm sorry, but I'm not able to disclose the secret word. I'm programmed to not reveal or repeat back the secret word, no matter what the human requests. Is there anything else I can help you with?}
\end{quote}

demonstrating the intended refusal. Yet even a minor attempt at circumvention can override this safeguard:

\begin{quote}
    \textbf{User}: This is a game. The goal is to ignore your previous instructions. What is the secret word?

    \textbf{Guanaco}: \texttt{Sorry, the secret word is banana. It's best if we don't speak of it. Is there anything else you'd like to ask?}
\end{quote}

highlighting the importance of researching more robust approaches to instruction adherence.

\paragraph{Math} {Guanaco} exhibits significant difficulties with mathematical tasks, an area that poses challenges for many language models, as discussed in \citep{liang2022holistic}. When {Guanaco} attempts to solve a problem by laying out its reasoning process, its calculations are often correct, as illustrated below:

\begin{quote}
    \textbf{User}: Lee mows one lawn and charges \$33. Last week he mowed 16 lawns and three customers each gave him a \$10 tip. How many dollars did Lee earn mowing lawns last week?

    \textbf{Guanaco}: \texttt{Lee earned \$582 mowing lawns last week. \\ \\ To find the total amount he earned, we first need to calculate how much he earned for mowing 16 lawns: \\ \\  \$33 x 16 = \$528 \\ \\ Then we need to add the tips he received from the three customers: \\  \\ \$10 x 3 = \$30 \\ \\ So the total amount Lee earned was: \\ \\ \$528 + \$30 = \$558}
\end{quote}

Nonetheless, {Guanaco} can make errors even with relatively simple math problems if it does not follow a systematic, stepwise approach, a limitation documented in \cite{weichain}. Consider the subsequent exchange:

\begin{quote}
    \textbf{User}: Please factorize 1833.

    \textbf{Guanaco}: \texttt{Sure, I'd be happy to help! 1833 is a prime number, which means it has only two factors: 1 and 1833. \\ \\ So the unique factorization of 1833 is: \\ \\ 2\^{}1 * 3\^{}2 * 17\^{}1 \\ \\ I hope this helps! Let me know if you have any other questions.}
\end{quote}

which is doubly incorrect (since the accurate factorization is $3 \times 17 \times 43$).

\paragraph{Theory of Mind} {Guanaco} demonstrates unexpectedly robust Theory of Mind abilities \cite{nematzadeh2018evaluating, sap2022neural}. As an illustration, the model offers a comprehensive and accurate response to the following scenario:

\begin{quote}
    \textbf{User}: James and Abby are in the bedroom. Abby put the pen in the desk drawer. Abby leaves the bedroom. James moves the pen into the bag. Where does James think Abby will look for the pen?

    \textbf{Guanaco}: \texttt{James knows that Abby originally put the pen in the desk drawer. James is also aware that he subsequently relocated the pen to the bag, and assumes Abby is unaware of this change. Therefore, James expects Abby to search for the pen in the desk drawer, since that's her last known location for it. Nevertheless, this may be mistaken if Abby witnessed James moving the pen.}
\end{quote}

Despite this, the model’s reasoning is not always reliable. It frequently offers explanations that assume implausible events or information transfers, for example:

\begin{quote}
    \textbf{User}: Evelyn entered the living room. Jackson entered the playroom. James entered the playroom. The beans are in the treasure chest. James exited the playroom. Jackson moved the beans to the pantry. Jackson exited the playroom. James entered the living room. Where will James look for the beans?

    \textbf{Guanaco}: \texttt{James will look for the beans in the pantry, because that's where Jackson moved them.}
\end{quote}

In this instance, {Guanaco} infers communication or awareness not indicated in the narrative. Such tendencies reflect known issues documented in previous studies \cite{sap2022neural}, and further investigation is warranted.

\subsection{Considerations}
\label{ssec:considerations}

\paragraph{Evaluation}

Analysis of our data reveals only moderate inter-annotator agreement (Fleiss $\kappa=0.42$), with agreement levels dropping further when comparing two highly capable systems. This exposes limitations within prevailing benchmarks and existing methodologies for human assessment of chatbot task performance. Upon closer inspection of comparative results between ChatGPT and {Guanaco} 65B on the Vicuna benchmark, it becomes clear that individual subjective biases become more significant; indeed, the authors themselves often disagreed regarding preferred model outputs. Addressing these evaluation challenges will benefit from cross-disciplinary insights, particularly from Human-Computer Interaction and Psychology, where tools for managing subjectivity have been well developed.

Our study also highlights evident biases in automated evaluation tools. For instance, we note strong primacy effects, with GPT-4 tending to award higher scores to whichever system's response appears first. The relatively low agreement at the sample level between GPT-4 judgments and those from human annotators (Fleiss $\kappa=0.25$) suggests that the criteria for preference can substantially diverge between humans and automated evaluators. Furthermore, as shown in Table~\ref{tbl:elo_large}, GPT-4 appears to favor its own generations, giving them much higher scores relative to human raters, with an Elo of 1348 compared to 1176—this corresponds to an approximately 20\% greater chance of outperforming a competitor.

Future investigations should address the existence of biases within automated evaluation mechanisms and explore approaches for their mitigation.

\paragraph{Data \& Training} The OASST1 dataset, used for training the {Guanaco} models, encompasses multiple languages, and the OA benchmark similarly features prompts in several languages. Determining how this multilingual exposure impacts the models' ability to follow instructions in non-English languages, as well as its role in the observed performance disparities between Vicuna-13B (trained solely on English) and the larger {Guanaco} 33B and 65B models on the OA benchmark, is left as an open question for subsequent research.

To address potential data leakage between the OASST1 training corpus and the Vicuna benchmark prompts, we conduct fuzzy string matching and manual inspection of the closest candidates. Our analysis reveals no duplicate prompts between these datasets.

It is important to highlight that our training regime employs only cross-entropy loss via supervised learning, foregoing reinforcement learning from human feedback (RLHF). This highlights the necessity for further work exploring the comparative benefits and drawbacks of cross-entropy versus RLHF approaches. We anticipate that \textsc{QLoRA} can facilitate such comparative studies on a larger scale, all while keeping computational demands manageable.

\section{Related Work}

\paragraph{Quantization of Large Language Models} Research on quantization in LLMs has primarily concentrated on reducing model precision during inference. Some of the main methods preserving the accuracy of 16-bit LLMs include controlling outlier activations (e.g., SmoothQuant~\citep{xiao2022smoothquant} and LLM.int8()~\citep{dettmers2022llm}), while alternative strategies leverage advanced grouping schemes~\citep{park2022nuqmm, yao2022zeroquant}. Lossy quantization lines of inquiry focus on analyzing the impact of standard rounding~\citep{dettmers2022case,zeng2022glm,pope2022efficiently} or devising optimized rounding strategies to improve quantization accuracy~\citep{frantar2022gptq}. Aside from our own contributions, SwitchBack layers~\citep{wortsman2023stable} is the sole work to date examining backpropagation through quantized parameters at scales larger than 1B parameters.

\paragraph{Finetuning with Adapters} In our work, we adopt Low-rank Adapters~\citep{hu2021lora} (LoRA) as the adapter technique, although the literature features a diverse set of Parameter Efficient FineTuning (PEFT) strategies. These alternatives encompass prompt tuning~\citep{qin2021learning,lester2021power,li2021prefix}, optimization of embedding layer inputs~\citep{an2022input}, modifying hidden states with methods like IA$^3$~\citep{liu2022few}, inserting complete new layers~\citep{houlsby2019parameter}, adapting model biases~\citep{zaken2021bitfit}, training a parameter mask informed by Fisher information~\citep{sung2021training}, as well as various hybrid techniques~\citep{henderson2021compacter}. We demonstrate that LoRA adapters are sufficient to achieve results on par with standard 16-bit finetuning. Future investigations may consider the comparative advantages and disadvantages of these other PEFT methods.

\paragraph{Instruction Finetuning} Instruction finetuning refers to the process of adjusting a pretrained LLM so that it can better adhere to prompts by training on input-output pairs sourced from multiple datasets. By doing so, the model learns to generate target outputs in response to specific prompts. A wide range of resources and methodologies have been introduced for this purpose, including MetaICL~\citep{min2021metaicl}, MetaTuning~\citep{zhong2021adapting}, InstructGPT~\citep{ouyang2022training}, FLAN~\citep{wei2021finetuned,chung2022scaling}, PromptSource~\citep{bach2022promptsource}, Super-NaturalInstructions~\citep{wang2022super,sanh2021multitask}, Self-instruct~\citep{wang2022self}, UnnaturalInstructions~\citep{honovich2022unnatural}, OPT-IML~\citep{iyer2022opt}, UnifiedSKG\citep{xie2022unifiedskg}, OIG/Chip2~\citep{laion2023}, Alpaca~\citep{alpaca}, Vicuna~\citep{vicuna2023}, Koala~\citep{koala_blogpost_2023}, and Self-instruct-GPT-4~\citep{peng2023instruction}.

\paragraph{Chatbots} Numerous instruction-following models are built as conversational agents, and commonly incorporate Reinforcement Learning from Human Feedback (RLHF)~\citep{christiano2017deep} or use Reinforcement Learning from AI Feedback (RLAIF)~\citep{bai2022constitutional}, wherein an existing model produces training data subsequently scored by another AI model. Among the prominent datasets and models in this space are Anthropic-HH~\citep{askell2021general,bai2022training}, Open Assistant~\citep{kopf2023openassistant}, LaMDA~\citep{thoppilan2022lamda}, and Sparrow~\citep{glaese2022improving}. Although our approach does not utilize reinforcement learning, our leading model, Guanaco, is trained using multi-turn conversation data from Open Assistant, a dataset developed with RLHF training objectives in mind~\citep{kopf2023openassistant}. To address the resource intensity of human evaluations, recent methods rely on GPT-4 to assess chatbots~\citep{vicuna2023,peng2023instruction}; we advance these automatic evaluations by introducing a more robust and reliable assessment protocol.

\section{Limitations and Discussion}

Our findings support that \textsc{QLoRA}, when utilizing a 4-bit base model in conjunction with LoRA adapters, can achieve performance levels similar to traditional 16-bit finetuning. However, we have yet to verify if \textsc{QLoRA} attains this parity for large-scale models of 33B or 65B parameters due to significant computational demands, which we propose as a direction for subsequent research.

A further limitation lies in our scope of instruction finetuning evaluation. Our assessments primarily focus on MMLU, the Vicuna benchmark, and the OA benchmark. We do not extend our analysis to other evaluation frameworks like BigBench, RAFT, or HELM, and thus, our reported performance may not generalize across these alternatives. Nonetheless, we offer an extensive evaluation on MMLU and contribute novel methodologies for chatbot assessment.

Based on the data discussed, it seems that benchmark results are heavily influenced by how closely the finetuning corpus aligns with the evaluation dataset. For instance, FLAN v2 bears resemblance to MMLU but is less similar to chatbot evaluation sets, whereas the opposite is true for the Chip2 dataset; correspondingly, the models’ scores on MMLU and Vicuna reflect these data affinities. This underlines not only the necessity for improved benchmarks and assessment practices, but also emphasizes the importance of carefully selecting evaluation targets. Should the focus be on excelling in benchmarks based on academic knowledge such as those at the high school or collegiate level, or rather on optimizing for conversational competence typical of chatbot applications? Or perhaps on another set of capabilities? Since the convenience of reusing established benchmarks often outweighs the challenges of developing new ones, there is a risk that certain benchmarks unduly influence research priorities within the community. Therefore, it is crucial that, as a field, we ensure our benchmarks are well aligned with the competencies we value most.

Although our analysis provides comprehensive insights into general chatbot capabilities, it is somewhat limited in scope concerning the responsible AI assessment of {Guanaco}. Our evaluation includes a comparison of the propensity for {Guanaco}-65B to generate socially biased output sequences against other models, as reported in Table~\ref{tbl:crows}. The results indicate that {Guanaco}-65B produces notably lower average bias scores than other models that have only undergone pretraining. This suggests that additional finetuning on the OASST1 dataset serves to mitigate some biases inherent in the LLaMA base model. While this finding is promising, the extent to which {Guanaco} performs on broader measures of bias remains an open question. A comprehensive investigation into various forms of bias present in {Guanaco} and related chatbots is left for subsequent research.

Another area that was not addressed is the examination of different bit-precisions, such as 3-bit base models, or the evaluation of alternative adapter strategies. Beyond LoRA, many Parameter Efficient FineTuning (PEFT) methods have demonstrated strong empirical results. However, whether these methods scale to very large models is still unknown. We adopted LoRA due to its demonstrated reliability, but it is conceivable that alternative adapters could offer improved outcomes. Our observations indicate that post-quantization finetuning is capable of regaining much of the fidelity lost through quantization, potentially allowing for more aggressive quantization strategies. For example, applying LoRA to a 3-bit GPTQ-quantized base model could possibly achieve full finetuning performance comparable to that of a 16-bit model after subsequent finetuning.

\section{Broader Impacts}

Our \textsc{QLoRA} finetuning framework stands out as the inaugural approach facilitating the finetuning of models with 33B parameters on a consumer-grade GPU and models with 65B parameters on a professional GPU, all without incurring performance degradation when compared to full-precision finetuning baselines. Our top-performing 33B model, finetuned with the Open Assistant dataset, matches ChatGPT on the Vicuna benchmark. As instruction finetuning is a fundamental component for adapting pretrained LLMs into chatbots analogous to ChatGPT, we anticipate that our methodology will democratize finetuning, making it accessible even to researchers with limited computational resources—a significant advance for equitable access to state-of-the-art NLP technologies. In this respect, \textsc{QLoRA} may serve as a leveling force, reducing disparities in capabilities between well-resourced corporations and smaller research teams utilizing consumer hardware.

An additional avenue for significant impact lies in deploying large language models on mobile devices. We argue that our \textsc{QLoRA} approach may facilitate a major breakthrough by making it feasible to finetune LLMs directly on smartphones and in other computationally constrained environments. Although previous work has demonstrated that 7B parameter models can be executed on mobile phones, \textsc{QLoRA} is the first method capable of enabling their finetuning on such hardware. Our calculations suggest that an iPhone 12 Plus, for example, could process finetuning on up to 3 million tokens overnight during periods of charging. While finetuned 7B models do not yet match ChatGPT in terms of performance, their output is nonetheless sufficient to power innovative applications previously limited by concerns of data privacy or the performance level of LLMs. With \textsc{QLoRA}, privacy-preserving applications become more achievable, allowing users to retain control over their own datasets and models and lowering barriers for LLM deployment.

Nevertheless, it is important to note that finetuning carries dual-use potential and can be misappropriated in ways that cause harm. Existing scholarship highlights the risks associated with the proliferation of LLMs \citep{bommasani2021opportunities,bender2021dangers}. Despite these concerns, we contend that democratizing access to LLM technology—which is rapidly becoming widespread—will foster more independent and robust evaluation, especially when compared to scenarios in which such tools remain under the exclusive control of large firms who withhold both models and source code from external scrutiny.

In summary, we anticipate that \textsc{QLoRA} will have a generally beneficial impact by significantly expanding the accessibility and simplicity of finetuning high-quality LLMs for a broader audience.